% --- Archon physics formalization source begin ---
% archon:physics
% archon:covers PhyXMiniProblems/problem_phyx_mini_0109.lean
% archon:source-report reports/phyx_mini/problem_phyx_mini_0109.source.json

\chapter{Physics problem phyx\_mini\_0109}
\label{ch:PhyXMiniProblems_problem_phyx_mini_0109}

\paragraph{Problem source.}
\#\# Physical scenario

Figure shows a small plant near a thin lens. The ray shown is one of the principal rays for the lens. Each square is 2.0 cm along the horizontal direction, but the vertical direction is not to the same scale. Use information from the diagram for the following.

\#\# Question

What is the focal length of the lens?

\#\# Answer choices

- A: A: 18.0cm

- B: B: 23.4cm

- C: C: 19.2cm

- D: D: 16.9cm

\#\# Figure caption (auxiliary; use the image as primary evidence)

The image shows a grid with a diagram illustrating the path of light rays interacting with a lens. Here’s a breakdown of the content:

- **Scenes and Objects:**
  - A **plant** is placed on the left side, labeled as "Plant."
  - A **lens** is represented in the center of the image by a blue, dashed rectangle labeled "Lens."
  - An **optic axis** is a horizontal line running through the center, labeled "Optic axis."

- **Light Rays:**
  - Two light rays are depicted:
    - A **solid purple ray** starts from the top of the plant, intersects the lens, and continues on a straight path rightward after refraction.
    - A **dashed purple ray** also starts from the top of the plant but reaches the lens along a different angle before being refracted to join the path of the solid ray past the lens.

- **Text and Labels:**
  - "Plant," "Lens," and "Optic axis" label their respective parts.
  - Question marks (?), in blue, are positioned vertically within the lens to possibly indicate unknown angles or refractive properties.

The diagram likely illustrates how light is refracted through a lens, such as in a basic optics experiment.

\#\# Dataset metadata

Geometrical Optics; Physical Model Grounding Reasoning, Spatial Relation Reasoning

\paragraph{Recorded answer/context.}
A: A: 18.0cm

\paragraph{Figure/image path.}
/root/proposal\_for\_physic/science-mango/phyx\_mini\_run/phyx\_data/test\_image/109.png

\paragraph{Formalization target.}
create a compiling Lean file with sorry bodies at `PhyXMiniProblems/problem\_phyx\_mini\_0109.lean`. The Lean declarations must preserve the physical quantities, dimensions or dimensional roles, figure labels, governing-law hypotheses, and final relation expressed by this problem.
Use Mathlib/Physlib names found through LeanExplore where available. If a physics API is missing, introduce faithful local abstractions rather than scalar placeholder aliases.

\begin{theorem}[Physics formalization target]
\label{thm:physics:phyx_mini_0109:target}
\lean{PhyXMiniProblems.ProblemPhyXMini0109.problem_phyx_mini_0109}
\uses{def:physics:phyx-mini-0109:phyxminiproblems-problemphyxmini0109-lengthincentimeters, def:physics:phyx-mini-0109:phyxminiproblems-problemphyxmini0109-principalraylensdiagram, def:physics:phyx-mini-0109:phyxminiproblems-problemphyxmini0109-matchesprimaryfigure, def:physics:phyx-mini-0109:phyxminiproblems-problemphyxmini0109-hascalibratedhorizontalgrid, def:physics:phyx-mini-0109:phyxminiproblems-problemphyxmini0109-hasphysicallengthmagnitudes, def:physics:phyx-mini-0109:phyxminiproblems-problemphyxmini0109-obeysthinlensprincipalraylaw, def:physics:phyx-mini-0109:phyxminiproblems-problemphyxmini0109-matchesanswerexactly}
The assigned autoformalize agent should translate this physics problem into Lean declarations in the covered file, with theorem and lemma proof bodies written as `by sorry`.
\end{theorem}
\begin{proof}
This is an autoformalization task, not a proof task. Produce statements that can later be proved by the physics prover without weakening the physical model.
\end{proof}

% --- Archon declaration topology begin ---
\section{Declaration topology}
\begin{definition}[Length Quantity]
  \label{def:physics:phyx-mini-0109:phyxminiproblems-problemphyxmini0109-lengthquantity}
  \lean{PhyXMiniProblems.ProblemPhyXMini0109.LengthQuantity}
  A signed physical length whose readout transforms coherently with units
\end{definition}
\begin{proof}
  This is the declared model definition of Length Quantity.
\end{proof}
\begin{definition}[centimeter Unit Choices]
  \label{def:physics:phyx-mini-0109:phyxminiproblems-problemphyxmini0109-centimeterunitchoices}
  \lean{PhyXMiniProblems.ProblemPhyXMini0109.centimeterUnitChoices}
  SI units with centimeters selected for length readouts
\end{definition}
\begin{proof}
  This is the declared model definition of centimeter Unit Choices.
\end{proof}
\begin{definition}[length In Centimeters]
  \label{def:physics:phyx-mini-0109:phyxminiproblems-problemphyxmini0109-lengthincentimeters}
  \lean{PhyXMiniProblems.ProblemPhyXMini0109.lengthInCentimeters}
  \uses{def:physics:phyx-mini-0109:phyxminiproblems-problemphyxmini0109-lengthquantity, def:physics:phyx-mini-0109:phyxminiproblems-problemphyxmini0109-centimeterunitchoices}
  The signed scalar readout of a physical length in centimeters
\end{definition}
\begin{proof}
  This is the declared model definition of length In Centimeters.
\end{proof}
\begin{definition}[Lens Approximation]
  \label{def:physics:phyx-mini-0109:phyxminiproblems-problemphyxmini0109-lensapproximation}
  \lean{PhyXMiniProblems.ProblemPhyXMini0109.LensApproximation}
  The thin-lens approximation stated in the problem
\end{definition}
\begin{proof}
  This is the declared model definition of Lens Approximation.
\end{proof}
\begin{definition}[Optic Axis Orientation]
  \label{def:physics:phyx-mini-0109:phyxminiproblems-problemphyxmini0109-opticaxisorientation}
  \lean{PhyXMiniProblems.ProblemPhyXMini0109.OpticAxisOrientation}
  Orientation of the labelled optic axis in the source figure
\end{definition}
\begin{proof}
  This is the declared model definition of Optic Axis Orientation.
\end{proof}
\begin{definition}[Vertical Scale Status]
  \label{def:physics:phyx-mini-0109:phyxminiproblems-problemphyxmini0109-verticalscalestatus}
  \lean{PhyXMiniProblems.ProblemPhyXMini0109.VerticalScaleStatus}
  Whether vertical page distances share the calibrated horizontal scale
\end{definition}
\begin{proof}
  This is the declared model definition of Vertical Scale Status.
\end{proof}
\begin{definition}[Figure Point]
  \label{def:physics:phyx-mini-0109:phyxminiproblems-problemphyxmini0109-figurepoint}
  \lean{PhyXMiniProblems.ProblemPhyXMini0109.FigurePoint}
  Points whose horizontal locations or optical roles occur in the diagram
\end{definition}
\begin{proof}
  This is the declared model definition of Figure Point.
\end{proof}
\begin{definition}[Figure Ray Part]
  \label{def:physics:phyx-mini-0109:phyxminiproblems-problemphyxmini0109-figureraypart}
  \lean{PhyXMiniProblems.ProblemPhyXMini0109.FigureRayPart}
  The three visually distinguished parts of the single depicted ray
\end{definition}
\begin{proof}
  This is the declared model definition of Figure Ray Part.
\end{proof}
\begin{definition}[Principal Ray Lens Diagram]
  \label{def:physics:phyx-mini-0109:phyxminiproblems-problemphyxmini0109-principalraylensdiagram}
  \lean{PhyXMiniProblems.ProblemPhyXMini0109.PrincipalRayLensDiagram}
  \uses{def:physics:phyx-mini-0109:phyxminiproblems-problemphyxmini0109-lengthquantity, def:physics:phyx-mini-0109:phyxminiproblems-problemphyxmini0109-lensapproximation, def:physics:phyx-mini-0109:phyxminiproblems-problemphyxmini0109-opticaxisorientation, def:physics:phyx-mini-0109:phyxminiproblems-problemphyxmini0109-verticalscalestatus, def:physics:phyx-mini-0109:phyxminiproblems-problemphyxmini0109-figurepoint, def:physics:phyx-mini-0109:phyxminiproblems-problemphyxmini0109-figureraypart}
  The physical quantities and incidence relations represented by the figure. axialPosition is only a horizontal coordinate. The plant height is retained as a physical quantity even though the diagram supplies no calibrated vertical readout. The relational ray fields preserve the fact that the dashed and solid incident pieces are one line without assigning page-picture coordinates physical meaning
\end{definition}
\begin{proof}
  This is the declared model definition of Principal Ray Lens Diagram.
\end{proof}
\begin{definition}[axial Separation In Centimeters]
  \label{def:physics:phyx-mini-0109:phyxminiproblems-problemphyxmini0109-axialseparationincentimeters}
  \lean{PhyXMiniProblems.ProblemPhyXMini0109.axialSeparationInCentimeters}
  \uses{def:physics:phyx-mini-0109:phyxminiproblems-problemphyxmini0109-lengthincentimeters, def:physics:phyx-mini-0109:phyxminiproblems-problemphyxmini0109-figurepoint, def:physics:phyx-mini-0109:phyxminiproblems-problemphyxmini0109-principalraylensdiagram}
  Signed horizontal separation from start to finish, in centimeters
\end{definition}
\begin{proof}
  This is the declared model definition of axial Separation In Centimeters.
\end{proof}
\begin{definition}[Matches Primary Figure]
  \label{def:physics:phyx-mini-0109:phyxminiproblems-problemphyxmini0109-matchesprimaryfigure}
  \lean{PhyXMiniProblems.ProblemPhyXMini0109.MatchesPrimaryFigure}
  \uses{def:physics:phyx-mini-0109:phyxminiproblems-problemphyxmini0109-lengthincentimeters, def:physics:phyx-mini-0109:phyxminiproblems-problemphyxmini0109-principalraylensdiagram}
  Primary-image and problem-text readouts. In particular, the backward-axis intercept and lens center are nine calibrated horizontal squares apart. This does not identify that separation with the unknown focal length
\end{definition}
\begin{proof}
  This is the declared model definition of Matches Primary Figure.
\end{proof}
\begin{definition}[Has Calibrated Horizontal Grid]
  \label{def:physics:phyx-mini-0109:phyxminiproblems-problemphyxmini0109-hascalibratedhorizontalgrid}
  \lean{PhyXMiniProblems.ProblemPhyXMini0109.HasCalibratedHorizontalGrid}
  \uses{def:physics:phyx-mini-0109:phyxminiproblems-problemphyxmini0109-lengthincentimeters, def:physics:phyx-mini-0109:phyxminiproblems-problemphyxmini0109-figurepoint, def:physics:phyx-mini-0109:phyxminiproblems-problemphyxmini0109-principalraylensdiagram, def:physics:phyx-mini-0109:phyxminiproblems-problemphyxmini0109-axialseparationincentimeters}
  The horizontal grid calibration: for two marked vertical grid lines, the physical axial separation is the counted number of squares times the physical width of one square. This is a general measurement rule, not an optical law
\end{definition}
\begin{proof}
  This is the declared model definition of Has Calibrated Horizontal Grid.
\end{proof}
\begin{definition}[Has Physical Length Magnitudes]
  \label{def:physics:phyx-mini-0109:phyxminiproblems-problemphyxmini0109-hasphysicallengthmagnitudes}
  \lean{PhyXMiniProblems.ProblemPhyXMini0109.HasPhysicalLengthMagnitudes}
  \uses{def:physics:phyx-mini-0109:phyxminiproblems-problemphyxmini0109-lengthincentimeters, def:physics:phyx-mini-0109:phyxminiproblems-problemphyxmini0109-principalraylensdiagram}
  Positivity of the three physical length magnitudes in the setup
\end{definition}
\begin{proof}
  This is the declared model definition of Has Physical Length Magnitudes.
\end{proof}
\begin{definition}[Obeys Thin Lens Principal Ray Law]
  \label{def:physics:phyx-mini-0109:phyxminiproblems-problemphyxmini0109-obeysthinlensprincipalraylaw}
  \lean{PhyXMiniProblems.ProblemPhyXMini0109.ObeysThinLensPrincipalRayLaw}
  \uses{def:physics:phyx-mini-0109:phyxminiproblems-problemphyxmini0109-lengthincentimeters, def:physics:phyx-mini-0109:phyxminiproblems-problemphyxmini0109-figurepoint, def:physics:phyx-mini-0109:phyxminiproblems-problemphyxmini0109-principalraylensdiagram, def:physics:phyx-mini-0109:phyxminiproblems-problemphyxmini0109-axialseparationincentimeters}
  The focal principal-ray law for a thin lens. If a non-axis-parallel principal incident ray has a backward extension meeting the optic axis and emerges parallel to that axis, then the axis intercept is the object-side principal focus. Consequently its axial distance from the lens optical center equals the focal-length magnitude. The law is quantified over candidate axis-intersection and lens-hit points; it contains neither the figure's grid count nor the requested numeric answer
\end{definition}
\begin{proof}
  This is the declared model definition of Obeys Thin Lens Principal Ray Law.
\end{proof}
\begin{definition}[Answer Choice]
  \label{def:physics:phyx-mini-0109:phyxminiproblems-problemphyxmini0109-answerchoice}
  \lean{PhyXMiniProblems.ProblemPhyXMini0109.AnswerChoice}
  The four displayed answer labels
\end{definition}
\begin{proof}
  This is the declared model definition of Answer Choice.
\end{proof}
\begin{definition}[answer Focal Length In Centimeters]
  \label{def:physics:phyx-mini-0109:phyxminiproblems-problemphyxmini0109-answerfocallengthincentimeters}
  \lean{PhyXMiniProblems.ProblemPhyXMini0109.answerFocalLengthInCentimeters}
  \uses{def:physics:phyx-mini-0109:phyxminiproblems-problemphyxmini0109-answerchoice}
  Focal-length readout printed beside each answer, in centimeters
\end{definition}
\begin{proof}
  This is the declared model definition of answer Focal Length In Centimeters.
\end{proof}
\begin{definition}[recorded Answer Choice]
  \label{def:physics:phyx-mini-0109:phyxminiproblems-problemphyxmini0109-recordedanswerchoice}
  \lean{PhyXMiniProblems.ProblemPhyXMini0109.recordedAnswerChoice}
  \uses{def:physics:phyx-mini-0109:phyxminiproblems-problemphyxmini0109-answerchoice}
  Dataset metadata: the recorded answer label, not a theorem premise
\end{definition}
\begin{proof}
  This is the declared model definition of recorded Answer Choice.
\end{proof}
\begin{definition}[Matches Answer Exactly]
  \label{def:physics:phyx-mini-0109:phyxminiproblems-problemphyxmini0109-matchesanswerexactly}
  \lean{PhyXMiniProblems.ProblemPhyXMini0109.MatchesAnswerExactly}
  \uses{def:physics:phyx-mini-0109:phyxminiproblems-problemphyxmini0109-lengthquantity, def:physics:phyx-mini-0109:phyxminiproblems-problemphyxmini0109-lengthincentimeters, def:physics:phyx-mini-0109:phyxminiproblems-problemphyxmini0109-answerchoice, def:physics:phyx-mini-0109:phyxminiproblems-problemphyxmini0109-answerfocallengthincentimeters}
  Exact agreement of a physical focal-length magnitude with one answer
\end{definition}
\begin{proof}
  This is the declared model definition of Matches Answer Exactly.
\end{proof}
\begin{lemma}[backward Axis Intercept Separation eq eighteen]
  \label{lem:physics:phyx-mini-0109:phyxminiproblems-problemphyxmini0109-backwardaxisinterceptseparation-eq-eighteen}
  \lean{PhyXMiniProblems.ProblemPhyXMini0109.backwardAxisInterceptSeparation_eq_eighteen}
  \uses{def:physics:phyx-mini-0109:phyxminiproblems-problemphyxmini0109-principalraylensdiagram, def:physics:phyx-mini-0109:phyxminiproblems-problemphyxmini0109-axialseparationincentimeters, def:physics:phyx-mini-0109:phyxminiproblems-problemphyxmini0109-matchesprimaryfigure, def:physics:phyx-mini-0109:phyxminiproblems-problemphyxmini0109-hascalibratedhorizontalgrid}
  The calibrated grid alone makes the horizontal separation between the backward intercept and lens center 9 * 2 = 18 cm. No optical interpretation of that separation is used in this intermediate statement
\end{lemma}
\begin{proof}
  The conclusion follows from the governing relations and figure readouts named in the statement of backward Axis Intercept Separation eq eighteen.
\end{proof}
% --- Archon declaration topology end ---
% --- Archon physics formalization source end ---
